\documentclass[10pt,a4paper]{amsart}
\usepackage[margin=19mm]{geometry}
\usepackage{amsmath,amssymb,mathtools}
\usepackage[scr=rsfs,cal=euler]{mathalfa}
\usepackage{xfrac}
\usepackage[unicode,hidelinks,bookmarks=false]{hyperref}
\newcommand{\ie}{i.e.\ }
\newcommand{\cf}{cf.\ }
\newcommand{\resp}{respectively\ }
\newcommand{\Subsection}{\subsection}
\providecommand{\nobreakdash}{\nobreak\hbox{-}}
\setlength{\emergencystretch}{2em}
\multlinegap0pt
\usepackage{bm}
\usepackage{amssymb}
\usepackage{enumerate}
\usepackage{xcolor}
\usepackage{tikz}
\usepackage[shortlabels]{enumitem}
\usepackage{float}
\newtheorem{theorem}{Theorem}
\newtheorem{claim}{Claim}
\newtheorem{poc}{Poc}
\newcommand*{\abs}[1]{\lvert#1\rvert}
\title{More on Nosal's spectral theorem: Books and $4$-cycles}
\author{Yongtao Li, Hong Liu, Shengtong Zhang}
\date{Source: arXiv:2508.14366v1 (CC BY 4.0)}
\begin{document}
\maketitle

\noindent\emph{Source and licence.} More on Nosal's spectral theorem: Books and $4$-cycles. Version arXiv:2508.14366v1 (CC BY 4.0), distributed by arXiv under the Creative Commons Attribution 4.0 International licence, http://creativecommons.org/licenses/by/4.0/, as recorded in the arXiv licence metadata for this exact version and shown on the abstract page https://arxiv.org/abs/2508.14366v1. Full licence notice: \texttt{licenses/arxiv-2508.14366-CC-BY-4.0.txt}.\par\smallskip

\noindent\textit{Editorial scope.} Counted unit: the article's theorem lem:Nosal-max-degree (source0.tex lines 305--310). The article's complete original proof of it is delivered with it (source0.tex lines 838--985, one \texttt{proof} environment of 148 lines, which the article places away from the statement). No mathematics is written, repaired or re-derived: the statement and the proof are the article's own lines, in the article's own notation, and no retained line is substituted or re-flowed.  No further source-local context is needed: every cross-reference of every retained line resolves inside the two delivered spans.  Nothing was omitted from any retained span, so the package carries no omission record.  No retained line contains a citation, so the wrapper carries no bibliography and no bibliography entry.  No retained line contains a figure environment, an image-inclusion command or drawing code, so no image is delivered and the package declares no asset dependency.  Editorial formatting only, in addition: an AMS \texttt{amsart} preamble that restates the article's own declarations and macros used by the retained lines, namely the environments \texttt{theorem}, \texttt{claim}, \texttt{poc} on separate counters, together with the standard packages those lines need.  The retained environments print in this document as Theorem 1, Claim 1, Poc 1 in order of appearance, whereas the article prints the same items with the numbering of its own full document.  The complete native source of the article as distributed in the arXiv e-print for this exact version is bundled unchanged as \texttt{sources/183545/source0.tex}.
\begin{theorem}
        \label{lem:Nosal-max-degree}
    If $G$ is an $m$-edge Nosal graph, 
    then for sufficiently large $m$,  
    \[ \Delta(G) \leq \frac{m}{2} + m^{0.99}.\] 
\end{theorem}

\begin{proof}[Proof of~Theorem \ref{lem:Nosal-max-degree}]
Let $i_*$ be the vertex that maximizes $x_{i_*}$. We repeat the following operation: if there is an edge $jk$ such that $j$ is not adjacent to $i_*$, then remove this edge from $G$ and add the edge $i_* j$. Note that this operation does not decrease $\bm{x}^T A_G \bm{x}$, so it preserves the Nosal property. When we can no longer do this operation, remove all isolated vertices, and let $G_*$ be the resulting graph. Then $G_*$ is Nosal, and $i_*$ is universal in $G_*$, that is, $i_*$ is adjacent to all the other vertices. Observe also that $\Delta(G_*) \geq \Delta(G)$. 

Thus, we may assume that our Nosal graph $G$ has a universal vertex $v$. Set $k = d_v$. For the sake of contradiction, 
we assume that $k > m/2 + m^{0.99}$.

Since $\lambda(G) > \sqrt{m}$, we have
$$ 2\sum_{\{i,j\} \in E} x_i x_j > \sqrt{m} 
\sum_{i \in V} x_i^2.$$
Let $G' = (V', E')$ be the subgraph of $G$ induced by 
$V(G) \backslash \{v\}$. 
Then $|V'|=k > m/2 +m^{0.99}$ and 
$|E'|=m-k$. 
We write $d_i'$ and $N'(i)$ for the degree and the neighborhood of a vertex $i\in V'$, respectively. 
Isolating the terms with $i = v$ in the above inequality, we have
$$2 \sum_{\{i,j\} \in E'} x_i x_j + 
2 \sum_{i \in V'} x_i x_v
> \sqrt{m} \sum_{i \in V'} x_i^2 + \sqrt{m} x_v^2.$$
Using the AM-GM inequality, we have
$$\sqrt{m} x_v^2 + \frac{1}{\sqrt{m}} 
\left( \sum_{i\in V'} x_i \right)^2 
\ge 2\sum_{i\in V'} x_i x_v.$$
Summing these two inequalities, we get 
\begin{equation} \label{eq-ge}
 2\sum_{\{i,j\} \in E'} x_i x_j + 
\frac{1}{\sqrt{m}} \left(\sum_{i \in V'} x_i \right)^2 
> \sqrt{m} \sum_{i \in V'} x_i^2. 
\end{equation} 
We now employ a classical estimate on 
$\sum_{\{i,j\} \in E'} x_i x_j$. Note that
$$ 2 \sum_{\{i,j\} \in E'} x_i x_j 
\leq \sum_{\{i,j\} \in E'} 
\left( \frac{\sqrt{d_j' + 1}}{\sqrt{d_i' + 1}} x_i^2 + \frac{\sqrt{d_i' + 1}}{\sqrt{d_j' + 1}} x_j^2\right).$$
For every $i$, we set
$$q_i := \sum_{j\in N'(i)} \frac{\sqrt{d_j' + 1}}{\sqrt{d_i' + 1}}.$$  
Then (\ref{eq-ge}) gives 
\begin{equation} \label{eq-recall}
\frac{1}{\sqrt{m}} \left(\sum_{i \in V'} x_i\right)^2 > 
\sum_{i \in V'} (\sqrt{m} - q_i ) \cdot x_i^2.
\end{equation}
We can bound $q_i$ as follows.

\begin{claim}\label{cl:1}
$q_i < \min\left\{ 
\sqrt{2(m - k)}, \sqrt{ \frac{d_i'}{d_i' + 1} \left((m - k) + d_i'^2\right)} \, \right\}.$   
\end{claim}


\begin{poc}
Using the Cauchy--Schwarz inequality, we obtain 
$$q_i \leq \sqrt{\frac{d_i'}{d_i' + 1}} \cdot \sqrt{\sum_{j \in N'(i)} (d_j' + 1)}.$$
Let $S$ denote the vertex set consisting of $i$ and its neighbors in $V'$. 
Then $\sum_{j \in N'(i)} (d_j' + 1)=\sum_{j\in S}d_j'$ is equal to twice the number of edges in $G'[S]$ plus the number of edges in $G'$ with exactly one vertex in $S$. In particular, this implies
$$\sum_{j \in N'(i)} (d_j' + 1) \leq 
\min\left\{2e(G'), e(G') + 
\binom{\abs{S}}{2}\right\}.$$
The claim follows as $e(G') = m - k$ and $\abs{S} = d_i' + 1$.  
\end{poc}

Recall that we assume $k > m/2 + m^{0.99}$. 
By Claim \ref{cl:1}, for each $i \in V'$ we have
$$\sqrt{m} - q_i > 0. $$ 
Applying the Cauchy--Schwarz inequality on the right hand side of~\eqref{eq-recall} gives 
$$\frac{1}{\sqrt{m}} \left(\sum_{i \in V'} x_i\right)^2>\sum_{i \in V'} (\sqrt{m} - q_i) \cdot x_i^2 \geq \left(\sum_{i \in V'} x_i \right)^2 \cdot \frac{1}{\sum\limits_{i \in V'} \frac{1}{\sqrt{m} - q_i}}.$$
So we must have
\begin{equation} \label{eqge-sqm}
    \sum_{i \in V'} \frac{1}{\sqrt{m} - q_i} > \sqrt{m}.
\end{equation}
Let $S_1$ be the set of vertices in $V'$ with degree at least $m^{0.25}$, and let $S_2 = V' \backslash S_1$.  
For any vertex $i \in S_1$, Claim \ref{cl:1} implies that  
\begin{equation}  \label{eq-S1} 
\frac{1}{\sqrt{m} - q_i} < \frac{1}{\sqrt{m} - \sqrt{2(m - k)}} < \frac{1}{\sqrt{m} - \sqrt{m - 2m^{0.99}}} < m^{-0.49}. 
\end{equation} 
For each vertex $i \in S_2$, by definition 
we have $d_i' < m^{0.25}$. 
Next, we prove a more sophisticated estimate 
on $1/(\sqrt{m} - q_i)$ for every $i\in S_2$. Let $d'_{i,1}$ denote the number of edges between $i$ and vertices in $S_1$, and $d'_{i,2}$ denote the number of edges between $i$ and vertices in $S_2$. 

\begin{claim}\label{cl:2}
    For any vertex $i\in S_2$, we have 
\begin{equation} \label{eq-S2}
\frac{1}{\sqrt{m} - q_i} \leq m^{-0.5} + (m^{-0.5} - m^{-0.6}) d'_{i,1} + m^{-0.6} d'_{i,2}.
\end{equation} 
\end{claim}


\begin{poc}
Using Cauchy-Schwarz's inequality 
as in Claim \ref{cl:1}, we have
$$\sum_{j \in S_1 \cap N'(i)} \frac{\sqrt{d_j' + 1}}{\sqrt{d_i' + 1}} \leq \sqrt{ \frac{d'_{i,1}}{d_i' + 1} \left((m - k) + d_i'^2\right)} \leq 
\sqrt{ \frac{d'_{i,1}}{d'_{i,1} + 1} 
\left( \frac{m}{2} - m^{0.98} \right)},$$
where the last inequality follows from 
$k> m/2 + m^{0.99}$ and $d_i'^2 < m^{0.5}$. 
Furthermore, we have
$$\sum_{j \in S_2\cap N'(i)} \frac{\sqrt{d_j' + 1}}{\sqrt{d_i' + 1}} \leq \sqrt{m^{0.25} + 1} \cdot \frac{d_{i,2}'}{\sqrt{d_i' + 1}} \leq m^{0.2} d_{i,2}'.$$
So for every $i\in S_2$, it follows that  
\[ q_i \leq \sqrt{ \frac{d'_{i,1}}{d'_{i,1} + 1} 
\left( \frac{m}{2} - m^{0.98} \right)} + m^{0.2} d_{i,2}' . \]
Then we have
$$\frac{1}{\sqrt{m} - q_i} \leq \frac{1}{\sqrt{m} - \sqrt{\frac{d'_{i,1}}{d'_{i,1} + 1} 
\left( \frac{m}{2} - m^{0.98} \right)} - m^{0.2} d_{i,2}'}.$$ 
Note that the difference between the first two terms in the denominator above is greater than $(1 - \sqrt{1 / 2})\sqrt{m}$, 
which together with $d_{i,2}' < d_i' < m^{0.25}$ 
yields 
$$\frac{1}{\sqrt{m} - q_i} \leq 
\frac{1}{\sqrt{m} - \sqrt{\frac{d'_{i,1}}{d'_{i,1} + 1} 
\left( \frac{m}{2} - m^{0.98}\right)}} + 100m^{-0.8} d_{i,2}'.$$
One can now directly verify
$$\frac{1}{\sqrt{m} - \sqrt{\frac{d'_{i,1}}{d'_{i,1} + 1} 
\left( \frac{m}{2} - m^{0.98}\right)}} \leq 
m^{-0.5} + (m^{-0.5} - m^{-0.6}) d'_{i,1}.$$
Indeed, if $d'_{i,1} \in \{0, 1\}$, then this follows by direct computation. If $d'_{i,1} \geq 2$ we have
$$\frac{1}{\sqrt{m} - \sqrt{\frac{d'_{i,1}}{d'_{i,1} + 1} \left( \frac{m}{2} - m^{0.98} \right)}} 
<  \frac{1}{\sqrt{m}} \cdot \frac{1}{1 - \sqrt{\frac{d'_{i,1}}{2(d'_{i,1} + 1)}}} 
< 0.9 m^{-0.5} (1 + d'_{i,1}),$$
where the last inequality follows since for any real number $x> 1.43$, 
\[ 0.9\left( 1-\sqrt{\frac{x}{2(x+1)}}\, \right)(1+x) >1. \]  
So we have verified the inequality (\ref{eq-S2}).  
\end{poc}

Combining (\ref{eq-S1}) with (\ref{eq-S2}), we have
\begin{align*}
 \sum_{i \in V'} \frac{1}{\sqrt{m} - q_i} &\leq m^{-0.49} \abs{S_1} + \sum_{i \in S_2} \left( 
m^{-0.5} + (m^{-0.5} - m^{-0.6}) d'_{i,1} + m^{-0.6} d'_{i,2} \right) \\
&= m^{-0.5} \abs{S_2} + m^{-0.49} \abs{S_1} + \sum_{i \in S_2} \left((m^{-0.5} - m^{-0.6}) d'_{i,1} + m^{-0.6} d'_{i,2} \right)
\end{align*}
We now bound the right hand side using the discharging method. For each edge $e \in E'$ and endpoint $i \in e$, we define a weight $w_{e, i}$ that $e$ ``discharges" to $i$ as follows
\begin{enumerate}
    \item If $e$ lies in $G'[S_1]$ or $G'[S_2]$, we set $w_{e, i} = m^{-0.6}$.
    \item Otherwise, let the endpoints of $e$ be $i, j$ with $i \in S_1$ and $j \in S_2$. We set $w_{e, i} = m^{-0.7}$ and $w_{e, j} = m^{-0.5} - m^{-0.6}$.
\end{enumerate}
Recall that $S_1$ denotes the set of vertices with degree at least $m^{0.25}$.  On the one hand, each vertex $i$ in $S_1$ satisfies
$$\sum_{e: i \in e} w_{e, i} \geq m^{-0.7} \cdot  m^{0.25} > m^{-0.49}$$
while each vertex $i$ in $S_2$ satisfies
$$\sum_{e: i \in e} w_{e, i} \geq (m^{-0.5} - m^{-0.6}) d'_{i,1} + m^{-0.6} d'_{i,2}.$$ 
On the other hand, each edge $e$ in $E'$ satisfies
\[ \sum_{i: i \in e} w_{e, i} \leq \max\{2m^{-0.6}, m^{-0.7} + m^{-0.5} - m^{-0.6}\} < m^{-0.5} - m^{-0.7} .\] 
So we conclude that
$$m^{-0.49} \abs{S_1} + \sum_{i \in S_2} (m^{-0.5} - m^{-0.6}) d'_{i,1} + m^{-0.6} d'_{i,2} 
\leq \sum_{i \in V', e \in E': i \in e} w_{e, i} \leq (m^{-0.5} - m^{-0.7}) \abs{E'}.$$
Invoking $|E'|=m-k$ and $|S_2| 
\le |V'| = k$, 
 we have
$$\sum_{i \in V'} \frac{1}{\sqrt{m} - q_i} \leq m^{-0.5} k + (m^{-0.5} - m^{-0.7})(m - k) \leq \sqrt{m},$$
which contradicts with (\ref{eqge-sqm}). 
Therefore, we must have $k \le m/2 + m^{0.99}$. 
\end{proof}

\end{document}
